\begin{abstract}
Single-cell trajectory inference from destructive time-course snapshots is fundamentally ill-posed: neither cross-time cell correspondences nor the continuous paths between snapshots are observed, so the observed snapshot distributions alone do not uniquely determine the underlying dynamics. 
Existing optimal transport and flow-based methods typically couple cells by Euclidean proximity at observed clock times, which can misalign trajectories when development is asynchronous and cells sampled at the same experimental time occupy different latent pseudotime stages. 
We propose PACE, a trajectory inference framework that selects geometry-consistent continuous transport dynamics from destructive time-course snapshots through three coupled components.
First, PACE constructs a state- and time-dependent anisotropic Riemannian metric that preserves low cost along locally supported tangent directions while penalizing normal velocity components.
Second, it alternates between refining cross-time couplings under the induced path-action cost and fitting endpoint-preserving neural bridges between adjacent snapshots.
Third, it distills the learned bridge dynamics into a global continuous-time velocity field over cellular states.
Across seven controlled and biological datasets covering nine held-out reconstruction experiments, PACE achieves the strongest overall reconstruction performance, reducing MMD, $\mathcal{W}_1$, and $\mathcal{W}_2$ by 23.7\% on average relative to the strongest competing baseline.
PACE also improves RNA-velocity alignment by 15.4\% on an embryoid body differentiation benchmark, without requiring explicit cell pairing, lineage tracing, or RNA velocity supervision during training.
% Anonymous-submission line (double-blind): keep for resubmission if needed.
% Code is available at \url{https://anonymous.4open.science/r/PACE-ICLR2027}.
Code is available at \url{https://github.com/AI4Science-WestlakeU/PACE}.
\end{abstract}
